\chapter{当代前沿研究简介}\label{ch:contemporary_research}

% 总论（不设节号；原「引言」与「主要方向总览」并入此处）
本书介绍了黎曼猜想知识谱系中的三个主要问题：
\begin{enumerate}
  \item \textbf{素数分布的主项}：素数定理所保证的主项渐近，以及显式公式中的主项
        （$\psi$ 侧的 $x$，$J$ 侧的 $\li(x)$ 等；已证）；
  \item \textbf{素数分布的精确结构}：显式公式给出计数函数与主项、
        \textbf{全体零点}（非平凡与平凡）及常数项等贡献的精确对应（结构已证）；
  \item \textbf{最优误差}：相对主项的最优误差阶，等价于非平凡零点均满足 $\operatorname{Re}\rho=1/2$，
        即黎曼猜想（未证）。
\end{enumerate}
当代工作大多仍落在第三项及其周边：如何证明（或定量逼近）全部非平凡零点在临界线上，
以及在尚未证明 RH 时，能推出多强的 $\psi$、$\pi$ 误差与零点统计。

\paragraph{黎曼猜想的证明难点}
由本书已建立的结果可知：显式公式给出素数计数与主项、全体零点及常数项等贡献的精确结构；
函数方程使 $\operatorname{Re}s=1/2$ 成为 $\xi$ 的对称轴；非平凡零点必在临界带
$0<\operatorname{Re}s<1$ 内。黎曼猜想所要求的，是进一步断言：
\textbf{每一个}非平凡零点均满足 $\operatorname{Re}\rho=1/2$。
该命题是全称陈述：任意虚部有限高度内的数值确认、以及临界线零点比例
$\kappa=\liminf N_0(T)/N(T)$ 的严格下界（即使 $\kappa>0$），均不能推出 RH；
而一个离线零点即可否证 RH。因此，证明必须在 $t\to\infty$ 与
$\operatorname{Re}s=1/2$ 的邻域上同时排除全部反例可能。

与此相对，已有若干\textbf{严格弱于} RH 的结果：素数定理依赖 $\zeta(1+it)\neq 0$；
自由区与密度定理给出 $\psi$、$\pi$ 的次优误差；亚凸界与临界线矩约束
$\zeta(\frac12+it)$ 的增长及均值。
对关联（Montgomery）描述临界线零点归一化间距的成对统计，
并猜想其与 \textbf{GUE}（Gaussian Unitary Ensemble，高斯酉系综——
随机 $N\times N$ 酉矩阵特征值的标准统计模型）一致；
详式见下文。对关联与数值验证增强对 RH 的信心，
但约束的是间距统计或有限高度，均不蕴含「全部零点在临界线上」。
Li 判据、Nyman--Beurling 条件、von Koch 型最优误差等与 RH \textbf{等价}，
只是同一断言的不同表述，并不拆成可逐个攻克的较弱中间命题。

从推理方向看：显式公式由零点导出算术误差；素数定理由 $\operatorname{Re}s=1$ 无零点导出主项。
RH 则要求临界带内零点的全称定位。现有方法长于边界排除与 $N(T)$ 计数，
短于「带内全部落在 $\operatorname{Re}s=1/2$」。
难点在于：精确结构与多种部分结果已知，全称定位尚未获得。
本章以下简介围绕该缺口的当代工作（比例、数值、误差、密度与谱设想）。

与本书章节的自然衔接：
\begin{itemize}
  \item 第~\ref{ch:zero_distribution}--\ref{ch:riemann_explicit_formula}~章：
        $N(T)$、显式公式、RH $\Leftrightarrow$ 最优误差阶（von Koch 型）；
  \item 第~\ref{ch:prime_number_theorem}~章与附录~\ref{app:pnt_estimates}：
        $\zeta(1+it)\neq 0$、自由区、密度型估计与 $\psi=x+o(x)$；
  \item 第~\ref{ch:riemann_numerical_approximation}--\ref{ch:visualization}~章：
        Hardy $Z$、Riemann--Siegel、$\pi_0$ 逼近与零点数据。
\end{itemize}

在\textbf{本书范围内}可优先了解的方向如下（按与主线的接近程度）：
\begin{enumerate}
  \item \textbf{临界线零点比例}：已证有多大比例的非平凡零点落在 $\operatorname{Re}s=1/2$ 上
        （Hardy--Selberg--Levinson--Conrey 线）；直接服务 RH 的「部分肯定」。
  \item \textbf{零点大规模数值确认}：有限高度内全部零点是否在临界线上
        （与第~\ref{ch:riemann_numerical_approximation}~章工具同一脉络）。
  \item \textbf{显式公式与误差阶}：RH $\Leftrightarrow$ $\pi-\Li$、$\psi-x$ 的最优型误差；
        无条件时由自由区 / 密度定理给出的次优误差。
  \item \textbf{零点间距与对关联}：Montgomery 猜想与 GUE；解释显式公式中振荡频率的统计。
  \item \textbf{亚凸界与 Lindelöf}：$\zeta(1/2+it)$ 的增长；弱于 RH 的增长控制，影响误差与均值。
  \item \textbf{零点密度定理}：$N(\sigma,T)$ 控制离线零点；通往 $\psi$、$\pi$ 的无条件误差。
  \item \textbf{谱理论 / 等价判据（简介）}：Hilbert--Pólya 设想，以及 Li、Nyman--Beurling、Lagarias 等
        与 RH 等价的分析表述（知其存在即可，本书不展开证明）。
\end{enumerate}
\textbf{本书不介绍}：椭圆曲线 BSD、自守形式与 Langlands、
$L$ 函数族的 Katz--Sarnak 全景、量子混沌的完整物理纲领、de Branges 路线等。
它们或属于更广的 $L$/算术几何，或与本书计数主线距离较远；
读者若有兴趣，可在学完 $\zeta$ 的延拓、函数方程与显式公式后再读专门文献。
下文各节均标明与上述三个问题及章节的对应；细节以原始文献为准。

% ============================================================
\section{临界线零点比例与数值确认}
% ============================================================

\subsection{临界线上零点的比例：从 Hardy 到 Conrey}

Hardy（1914）证明临界线 $\operatorname{Re}s=\tfrac{1}{2}$ 上有无穷多个零点，但未给出占全体非平凡零点的比例。
定量刻画该比例，是「向 RH 靠近」的经典解析课题。

\begin{definition}{临界线零点比例}{}
设 $N(T)$ 为虚部在 $(0,T]$ 内的非平凡零点总数（含重数；见第~\ref{ch:zero_distribution}~章），
$N_0(T)$ 为其中位于临界线上的个数。定义
\[
  \kappa = \liminf_{T\to\infty} \frac{N_0(T)}{N(T)}.
\]
黎曼猜想（RH）蕴含（且在标准约定下等价于）$\kappa = 1$。
\end{definition}

主要历史节点：

\begin{center}
\renewcommand{\arraystretch}{1.5}
\begin{tabular}{lll}
\toprule
\textbf{年份} & \textbf{作者} & \textbf{结论} \\
\midrule
1914 & G.\,H.\ Hardy & $N_0(T)\to\infty$ \\
1921 & Hardy \& Littlewood & $N_0(T) \gg T$ \\
1942 & A.\,Selberg & $N_0(T) \gg T\log T$（正比例） \\
1974 & N.\,Levinson & $\kappa \geq \tfrac{1}{3}$ \\
1989 & J.\,B.\ Conrey & $\kappa \geq \tfrac{2}{5}$ \\
2011 & Bui, Conrey \& Young & $\kappa \geq 0.4128$ \\
2013 & Feng & $\kappa \geq 0.4138$ \\
\bottomrule
\end{tabular}
\end{center}

Levinson--Conrey 路线的核心是将 $\zeta$ 与 $\zeta'$ 通过 Mollifier（磨光子）结合，
用平均值定理估计「同时接近零」的尺度。
这些估计依赖第~\ref{ch:zero_distribution}~章的 $N(T)$ 与第~\ref{ch:riemann_zeta_continuation}~章的函数方程，
是本书已备工具的精密综合。
\textbf{注意}：$\kappa>0$ 的下界\textbf{不是} RH；RH 要求 $\kappa=1$。
已证比例约四成，说明「多数零点在临界线上」有强证据，但与「全部」之间仍有本质差距。

\subsection{大规模数值确认}

与「比例定理」并行，计算方面用 Riemann--Siegel / Hardy $Z$
（第~\ref{ch:riemann_numerical_approximation}~章）在有限高度内逐个确认零点位置：

\begin{center}
\renewcommand{\arraystretch}{1.5}
\begin{tabular}{lll}
\toprule
\textbf{年份} & \textbf{作者 / 工具} & \textbf{已验证范围（摘要）} \\
\midrule
1987 & A.\,Odlyzko & $10^{20}$ 附近 $10^7$ 个零点的高精度统计 \\
2001 & van de Lune 等 & 前约 $1.5\times 10^9$ 个零点在临界线上 \\
2004 & Gourdon \& Demichel & 前 $10^{13}$ 个零点在临界线上 \\
2021 & Platt \& Trudgian & 前 $3\times 10^{12}$ 个零点的严格验证 \\
\bottomrule
\end{tabular}
\end{center}

在各自工作所达范围内，未发现离线零点。
须与第~\ref{ch:zero_distribution}~章一致理解：RH 是全称命题，\textbf{一个}离线零点即可否证；
上表是有限高度内的确认，\textbf{不是} RH 的证明。
其理论价值首先在于\textbf{可证伪性}（把可能的反例推到极高虚部），
其次为对关联、矩猜想、$\pi_0$ 逼近等提供零点数据
（见第~\ref{ch:riemann_numerical_approximation}--\ref{ch:visualization}~章）。

\subsection{计算与理论的互动}

\begin{enumerate}
  \item \textbf{统计猜想的数值支托}：Odlyzko 等高精度数据使 Montgomery 对关联与 GUE 的吻合成为
        「最强旁证」之一，并推动临界线矩的系统猜想（下节）。
  \item \textbf{反例边界}：若存在反例，其虚部必高于已严格检验的高度——为解析方法设立基准。
  \item \textbf{与本书 $\pi_0$ 逼近}：同一批 $\{\gamma_k\}$ 进入显式公式截断，
        在算术轴上拟合 $\pi$（第~\ref{ch:riemann_numerical_approximation}--\ref{ch:visualization}~章）；
        数值确认零点 $\Leftrightarrow$ 数值可用的振荡频率表。
\end{enumerate}

% ============================================================
\section{显式公式、最优误差与无条件误差}
% ============================================================

本节把当代结果扣回本书主线：\textbf{显式公式已给出通项结构；RH 管误差阶}。

\subsection{RH 与最优误差阶（温故）}

第~\ref{ch:riemann_explicit_formula}~章：在合适意义下
\[
\psi_0(x)=x-\sum_{\rho}\frac{x^{\rho}}{\rho}-\log(2\pi)-\frac12\log(1-x^{-2}),
\]
以及 $J$、$\pi$ 侧的等价形式。主项与零点项的\textbf{结构}已证；
RH（$\operatorname{Re}\rho=1/2$）与 von Koch 型最优误差等价，例如
\[
\pi(x)-\Li(x)=O\bigl(\sqrt{x}\,\log x\bigr)
\quad\Longleftrightarrow\quad
\operatorname{Re}\rho=\frac12
\]
（表述细节与 $\psi$ 侧 $O(\sqrt{x}\log^2 x)$ 等见该章）。
因此当代任何「证明 RH」的路线，在素数分布语言下都等于给出上述误差；
任何「部分结果」则对应\textbf{弱于最优}但仍有用的误差或零点约束。

\subsection{无条件误差从何而来}

在未设 RH 时，经典路线是：
\begin{enumerate}
  \item $\zeta(1+it)\neq 0$（第~\ref{ch:prime_number_theorem}~章）保证主项不被同阶破坏；
  \item \textbf{零点自由区}（附录~\ref{app:pnt_estimates}）把围道左移到 $\sigma=1-c/\log|t|$ 一类区域，
        得 $\psi(x)=x+O\bigl(x\exp(-c\sqrt{\log x})\bigr)$ 等 de la Vallée Poussin 型误差；
  \item \textbf{零点密度定理}（下节）控制 $\operatorname{Re}\rho>\sigma$ 的零点个数，
        用于短区间、算术级数及多种均值型估计。
\end{enumerate}
当代改进多在：扩大自由区、加强密度指数、或在更短区间上得到素数。
\textbf{这些结果全部落在「精确结构已有、最优误差未达」的中间地带}，
与上文总论的三个主要问题划分一致。

% ============================================================
\section{零点间距、对关联与临界线矩}
% ============================================================

\subsection{Montgomery 对关联猜想}

1972 年，Montgomery 在研究非平凡零点的对关联时提出：

\begin{definition}{Montgomery 对关联猜想}{}
设 $0 < \gamma_1 \leq \gamma_2 \leq \cdots$ 为 $\zeta$ 的非平凡零点虚部（常在 RH 下叙述，
并以 $\tfrac{2\pi}{\log T}$ 归一化）。对关联函数
\[
  R_2(\alpha) = \lim_{T\to\infty} \frac{1}{N(T)}
  \#\!\left\{(\gamma_m,\gamma_n):\, 0 < \gamma_m, \gamma_n \leq T,\, 0<\frac{(\gamma_m-\gamma_n)\log T}{2\pi}\leq\alpha\right\}
\]
被猜想等于
\[
  R_2(\alpha) = \int_0^\alpha \left[1 - \left(\frac{\sin\pi u}{\pi u}\right)^2\right]\mathrm{d}u,
\]
即与上文所引 GUE 特征值间距的对关联一致。
\end{definition}

\begin{explainframe}
\noindent\textbf{与本书显式公式的联系。}
显式公式中每一非平凡零点贡献一项以 $\gamma$ 为频率的振荡
（$\psi$ 侧 $x^{\rho}/\rho$，$J$ 侧 $\li(x^{\rho})$ 等；第~\ref{ch:riemann_explicit_formula}~章）。
对关联刻画的是这些频率两两靠近的统计，因而关系到误差项的涨落结构，
而不仅是单个零点的位置。
Montgomery--Dyson 与 GUE 的对照、以及 Odlyzko 的大规模计算，
是该统计图景的经典来源与数值支持（详见本节定义与下文数值确认）。
\end{explainframe}

\subsection{Keating--Snaith：临界线上 $|\zeta|$ 的矩}

2000 年，Keating 与 Snaith 用随机矩阵的特征多项式模型预测临界线矩：

\begin{definition}{$\zeta$ 的临界线矩}{}
对 $k>-\tfrac{1}{2}$，
\[
  I_k(T) = \frac{1}{T}\int_0^T \left|\zeta\!\left(\tfrac{1}{2}+it\right)\right|^{2k} \mathrm{d}t.
\]
猜想 $I_k(T) \sim c_k\,(\log T)^{k^2}$，其中 $c_k$ 由 GUE 特征多项式理论给出。
\end{definition}

已严格知道的情形包括：
$k=1$（Hardy--Littlewood，$I_1\sim\log T$）；
$k=2$（Ingham，$I_2\sim \tfrac{1}{2\pi^2}(\log T)^4$）；
更高整数 $k$ 多依赖 RH 或更强假设的条件结果。
这与「$\zeta$ 在临界线上的值分布」直接相关，并与第~\ref{ch:riemann_numerical_approximation}~章的数值图像一致；
\textbf{本书不}展开一般 $L$ 函数族的 Katz--Sarnak 分类（属更广纲领）。

% ============================================================
\section{亚凸界、Lindelöf 与零点密度}
% ============================================================

\subsection{临界带内的增长与 Lindelöf 猜想}

由 Phragmén--Lindelöf（第~\ref{ch:riemann_zeta_continuation}~章思路），临界带内有凸型界；
RH 蕴含更强的 Lindelöf 猜想：

\begin{definition}{Lindelöf 猜想}{}
\[
  \zeta\!\left(\tfrac{1}{2}+it\right) = O\!\left(|t|^\varepsilon\right),
  \quad \forall\,\varepsilon>0.
\]
Lindelöf 猜想是 RH 的推论，本身仍未证明。
\end{definition}

无条件亚凸界的当代代表结果之一是：

\begin{theorem}{Bourgain，2017}{}
\[
  \zeta\!\left(\tfrac{1}{2}+it\right) = O\!\left(|t|^{13/84+\varepsilon}\right).
\]
指数 $13/84 \approx 0.1548$ 改进了此前 Huxley 等结果。
\end{theorem}

亚凸界弱于 Lindelöf / RH，但对多种与 $\zeta$ 均值、某些素数问题相关的估计仍有直接用处；
其技术核心是指数和与 van der Corput 型方法，\textbf{不属于}本书证明提要范围，仅记结论与位置。

\subsection{零点密度定理}

记
\[
  N(\sigma, T) = \#\!\left\{\rho = \beta+i\gamma:\,\zeta(\rho)=0,\,
  \beta>\sigma,\,0<\gamma\leq T\right\}.
\]

\begin{theorem}{零点密度定理（Ingham 型，示意）}{}
对 $\tfrac{1}{2}\leq\sigma<1$，有
\[
  N(\sigma,T) \ll T^{A(1-\sigma)}\log^B T
\]
一类上界（经典 Ingham 取 $A=\tfrac{12}{5}$；后世有多种改进）。
\end{theorem}

$N(\sigma,T)$ 回答：「有多少零点落在 $\operatorname{Re}s>\sigma$？」
它弱于 RH（RH 要求 $\sigma>1/2$ 时 $N(\sigma,T)=0$ 对大 $T$），
但已能推出若干与 $\psi$、$\pi$ 相关的误差与均值结果（见 Iwaniec--Kowalski 等）。
与本书第~\ref{ch:prime_number_theorem}~章的分工：
该章给 $\psi\sim x$ 的逻辑骨架；密度与自由区的定量细节在附录与专门教本中。
\textbf{由密度导出的短区间素数、算术级数等}可视为同一工具箱的应用；
有界间隙、孪生等加性问题仍\textbf{不在}本书范围（见前言）。

% ============================================================
\section{谱理论途径与 RH 的等价表述（简介）}
% ============================================================

\subsection{Hilbert--Pólya 设想}

Hilbert 与 Pólya 各自猜想（约 1910--1920 年代，未形成当时的正式论文表述）：

\begin{definition}{Hilbert--Pólya 猜想}{}
存在自伴算子 $\mathcal{H}$，其特征值集合恰为非平凡零点的虚部 $\{\gamma_k\}$。
则 RH 对应「谱全为实」。
\end{definition}

Montgomery--Odlyzko 的 GUE 吻合，常被解读为：若这样的 $\mathcal{H}$ 存在，
应与量子混沌中的哈密顿量属同一统计类。该设想把 RH 改写为谱问题，
但\textbf{构造}出正确的 $\mathcal{H}$ 仍是开放的。

\subsection{Berry--Keating 与 Connes（提要）}

\paragraph{Berry--Keating（1999）}
建议经典 $H=xp$ 的量子化与零点谱相关；半直线上的 $\hat x\hat p+\hat p\hat x$ 等模型
至今仍是数学物理中的活跃试验场，\textbf{尚未}成为 RH 的证明。

\paragraph{Connes 非交换几何途径（1999）}
构造与 Adele 上迹公式相联系的框架，使迹公式在形式上对应 Weil 显式公式
——而 Weil 公式是第~\ref{ch:riemann_explicit_formula}~章黎曼显式公式的推广。
\begin{explainframe}
\noindent\textbf{显式公式作为共同语言。}
本书的 $\psi_0$ 显式公式把算术阶梯与零点级数写成恒等式。
Weil 将其推广到更一般的测试函数；Connes 等则试图为「零点与谱」的对应提供算子论表述。
谱理论途径与本书主线的接点，仍是显式公式（或其 Weil 型推广）所给出的
素数与零点之间的恒等结构。
\end{explainframe}

\subsection{若干等价判据（仅列名）}

下列条件均与 RH \textbf{等价}（在标准解析数论假设下），本书\textbf{不}推导：
\begin{itemize}
  \item \textbf{Li 判据}：某列实系数 $\lambda_n\geq 0$；
  \item \textbf{Lagarias 初等型不等式}：涉及调和数与 $\sigma(n)$ 的不等式对一切 $n$ 成立；
  \item \textbf{Nyman--Beurling}：某函数空间中 $1$ 可用 $\zeta$ 相关函数逼近的完备性表述。
\end{itemize}
它们说明：RH 有多种与「$\operatorname{Re}\rho=1/2$」等价的分析/初等包装；
学习上仍建议以零点几何与显式公式误差为主干，判据作为延伸阅读。

% ============================================================
\section{本书知识谱系与延伸阅读}
% ============================================================

\begin{center}
\renewcommand{\arraystretch}{1.7}
\begin{tabular}{|>{\raggedright\arraybackslash}p{4.6cm}|>{\raggedright\arraybackslash}p{4.4cm}|>{\raggedright\arraybackslash}p{5cm}|}
\hline
\textbf{方向（本书范围内）} & \textbf{对应章节} & \textbf{延伸阅读（示例）} \\
\hline
临界线零点比例 & 第\ref{ch:zero_distribution}章、$N(T)$ & Iwaniec--Kowalski，\textit{Analytic Number Theory} \\
\hline
数值确认与 $\pi_0$ 数据 & 第\ref{ch:riemann_numerical_approximation}--\ref{ch:visualization}章 & Odlyzko 零点表；Platt--Trudgian（2021） \\
\hline
显式公式与最优误差 & 第\ref{ch:riemann_explicit_formula}章、本章总论 & Edwards；Titchmarsh \\
\hline
对关联 / GUE & 第\ref{ch:riemann_explicit_formula}章（振荡频率） & Montgomery；Keating--Snaith（2000） \\
\hline
亚凸界、Lindelöf & 第\ref{ch:riemann_zeta_continuation}章（增长） & Bourgain（2017）及评述 \\
\hline
自由区、密度、$PNT$ 误差 & 第\ref{ch:prime_number_theorem}章、附录\ref{app:pnt_estimates} & Davenport；Titchmarsh \\
\hline
Hilbert--Pólya / 谱设想 & 显式公式；零点几何 & Berry--Keating，\textit{SIAM Review}（1999） \\
\hline
\end{tabular}
\end{center}

\begin{center}
\fbox{\parbox{0.92\textwidth}{
\centering
\vspace{0.2em}
\textbf{第~\ref{ch:contemporary_research}~章：与本书议题相关的进展速览}\\[6pt]
\renewcommand{\arraystretch}{1.6}
\begin{tabular}{lp{9.2cm}}
\hline
\textbf{方向} & \textbf{要点} \\
\hline
临界线比例 & Conrey $\kappa\geq 2/5$；Feng $\kappa\geq 0.4138$（$\neq$ RH） \\
数值确认 & 极高高度内无离线零点；可证伪，非证明 \\
最优误差 & RH $\Leftrightarrow$ von Koch 型 $\pi-\Li$ / $\psi-x$ 误差 \\
无条件误差 & 自由区 + 密度定理 $\Rightarrow$ 次优但有效的 $\psi$ 误差 \\
对关联 / 矩 & Montgomery--GUE；Keating--Snaith 临界线矩 \\
亚凸界 & Bourgain $|\zeta(1/2+it)|\ll |t|^{13/84+\varepsilon}$ \\
谱途径 & Hilbert--Pólya 等：开放；显式公式为共同语言 \\
\hline
\end{tabular}
\vspace{0.2em}
}}
\end{center}

\section*{本章小结与后续展望}
\addcontentsline{toc}{section}{本章小结与后续展望}

本章在总论所列范围内，简介了与零点、显式公式误差及数值确认相关的当代方向（见上文各节与速览框）。
若继续深入，可优先：（1）Titchmarsh / Edwards 中与零点、显式公式相关的章节；
（2）零点数值与对关联的经典论文；
（3）再视兴趣进入谱理论或更广的 $L$ 函数。
主体理论与三个主要问题的划分以第~\ref{ch:riemann_explicit_formula}~章及本章总论为准。
